\scanpage{64}
contient pas \uline{assez} de \inserted{$\Sigma$-}projectifs; il faut y rajouter $\Psi^0$\ldots) \deleted{Il est vrai aussi} Ces \deleted{\unsure{mêmes}} assertions sont d'ailleurs valables aussi dans la plus gde catégorie $P$ [détermination des \inserted{$\Sigma$-}proj. et des $\Sigma$-inj. et le fait qu'il y en ait assez].

Les $\operatorname{Ext}^{i}$ calculés par les résolutions injectives \inserted{à droite}, ou projectives \deleted{\unsure{sont (à gauche}} \inserted{à g.} ou les deux simultanément, sont les \unsure{mê.|$\mathrm{Ext}$} pour $i\geq1$ ce sont les hyperext habituels de \inserted{$L_{\bullet}$ et de $M_{\bullet}$}, pour $i=0$ ce sont par contre les $Z^i$.

Si $F:Q=P_1\hookrightarrow P\longrightarrow B$ ($k$-add.) est un foncteur, on peut construire les foncteurs dérivés \inserted{à droite ou gauche} en utilisant des résolutions injectives ou projectives. Nous sommes intéressés notamment aux $R^iF(\Psi^0)$, qu'on notera aussi $H^i(F)$ [on doit pouvoir montrer que ce sont les dérivés \inserted{droits} de $F\longmapsto F(\Psi^0)\overset{\mathrm{déf}}{=}\Gamma(F)$, quand $F$ varie dans la cat. des foncteurs \inserted{$k$-add} $P\longrightarrow \underline{B}$ \inserted{$\Sigma$-}exacts à gauche, (si on se restreint à de tels $F$); qu'en effet cette catégorie est équivalente: \deleted{\unsure{del.}} $\deleted{\unsure{B}} B^{\bullet}$, $F\longmapsto F(\Psi^0)$ s'identifie, \deleted{\unsure{alors}} foncteur $C^{\bullet}\longmapsto Z^0(C^{\bullet})\inserted{{}=H^0(C^{\bullet})}$ sur $B^{\bullet}$, dont les dérivés sont comme sont les $H^i(C^{\bullet})$ pour $i\geq1$. On a aussi $RF(\Psi^0)\simeq F(\mathbb{C}^{\bullet})$
% archive sheet number: 63
